\section{Related Work}\label{sec:background}
\subsection{Code Style Analysis}
In academia, the importance of code quality cannot be overstated. Although problem-solving and algorithm implementation often take precedence, the clarity and readability of the code are crucial for effective collaboration and maintenance in software development \cite{9637193}. Despite this, many programming courses do not adequately emphasize code quality \cite{Chen2018AnAA}. This is often because the time-consuming nature of manual review forces instructors to focus on assignments that assess problem-solving and algorithmic skills rather than coding standards \cite{9637193, 9766807}.

The inclusion of code analysis tools in academic settings, as explored in~\cite{10590588,10.1007/978-3-031-72483-1_1, Imhmed2024fosterCQ}, provides a systematic approach for gaining insights into student code quality and performance. Albluwi and Salter~\cite{albluwi2020using} highlighted the effectiveness of feedback from static analysis tools in reducing style errors over time, indicating that these tools can play a vital role in improving code quality education. These academic efforts to highlight and improve code quality through teaching strategies and static analysis tools are pivotal. They enhance students' coding skills and prepare them for software development's collaborative and evolving nature in academia and the professional world.

\subsection{Large Language Models in Test Case Generation}
LLM integration in programming education shows significant promise. LLMs can automate test case generation from software requirements or user stories, replacing the traditional manual approach, which is error-prone and labor-intensive. This automation improves turnaround times and reduces the need for human oversight~\cite{kumar2024using}. LLMs can also foster an interactive learning environment by delivering real-time feedback, accelerating students' iterative coding process, and reinforcing software testing principles~\cite{xiong2023program, dobslaw2025challenges}. Their scalability makes them particularly valuable in large programming courses, ensuring consistent and comprehensive test coverage across diverse projects. Furthermore, custom applications of LLMs support various assessment formats, facilitating innovative pedagogical approaches such as adaptive testing~\cite{xue2024llm4fin}. By catching errors early and improving software reliability, LLM-driven test case generation enhances student learning and aligns with best practices in modern software development.

\subsection{Automated Code Assessment in Education}
Several automated code assessment tools and techniques have been proposed for educational purposes. Cipriano et al.~\cite{cipriano2022drop} of Lusófona University proposed an automatic assessment tool used across several computer science courses. The proposed command-line tool leverages Apache Maven and Unit testing to automate software build and functionality testing. Buffardi \& Edwards~\cite{Buffardi2013} utilized Web-CAT to automate the evaluation of student code submissions for an introductory programming course at Virginia Tech, allowing students to repeatedly revise and re-submit their code for feedback on their code implementation and functionality. Heckman \& King~\cite{Heckman2018} of North Caloraina State University proposed a system that utilizes Jenkins and GitHub for code assessment.

Several studies~\cite{5557837,7883336,vander2013improving, peveler2017submitty} utilized unit testing to evaluate student code for correctness and functionality. Insa and Silva~\cite{INSA201859} provided black-box and white-box testing APIs for code assessment, along with a feature that allows students to compare the output of their code with the expected output. Peveler et al.~\cite{peveler2017submitty} used several assessment techniques, including the Linux diff tool for output comparison, checks for memory leaks, and performing static code analysis.